\section{Correctness Properties}\label{sec:correctness}

The \bxthree{} Framework is designed to satisfy a set of independently stated
correctness properties that together constitute the \textit{upstream accountability
guarantee}: that consequential decision authority in any \bxthree{} system
remains permanently attributable to a named human principal, and that any
unresolved exception condition is compulsorily escalated to that principal
without passing through any intermediate machine actor as a terminus.

Each property is stated as a formal invariant over the full framework state space,
supported by a proof sketch that traces the structural mechanisms enforcing it,
and is cross-referenced to the specific pillar or layer that provides the
enforcing mechanism.  The properties are not independent -- they are mutually
supporting, and each depends on the correct operation of specific \fact{} layer
pre-commit hooks that are identified explicitly.

% ---------------------------------------------------------------
\subsection{Property 1: Intent Preservation}\label{sec:intent-preservation}

\theorem[Intent Preservation]{Every consequential action executed by any node $n$ in any \bxthree{} system is authorized by a \fact{} layer entry in the authorization ledger and falls within the operating range $R(n)$ defined by the \purpose{} layer of $n$.}

\subparagraph{Formal statement.}  Let $\mathcal{A}(n)$ denote the set of all consequential actions executed by node $n$ during its operational lifetime.  Let $\mathcal{E}(n)$ denote the set of authorization ledger entries for node $n$.  Intent Preservation requires:
\[
\forall a \in \mathcal{A}(n)\ :\ \exists e \in \mathcal{E}(n)\ :\ \text{action}(e) = a
\ \land\ \text{scope}(e) \subseteq R(n)
\]
where $R(n)$ is the operating range defined by the \purpose{} layer at node provisioning.

\subparagraph{Enforcing mechanism.}  The \fact{} layer's pre-commit hook
intercepts every action request from the \bounds{} layer before it is dispatched
to the execution engine.  The hook consults the authorization ledger to verify
that a valid entry exists for the requested action within the current operating
range.  If no valid entry exists -- if the action is unauthorized or the scope
parameter of the existing entry does not cover the current execution context --
the pre-commit hook rejects the action and the node enters the Bailout Protocol
at the BOUNDED state.  No machine actor, including the \bounds{} layer, may
bypass this pre-commit check.

\subparagraph{Dependence.}  This property depends on Invariant~\ref{inv:intent}
(the per-node intent preservation invariant established at provisioning) and on
Invariant~\ref{inv:ledger-immutability} (the immutability of the authorization
ledger).  If the ledger were mutable, a compromised machine actor could retroactively
authorize an unauthorized action; if the intent preservation invariant were
violated, the operating range check at pre-commit would lack a valid reference
scope.

% ---------------------------------------------------------------
\subsection{Property 2: Structural Auditability}\label{sec:structural-auditability}

\theorem[Structural Auditability]{The complete action history, exception history, and resolution history of any node $n$ is recoverable from the Forensic Ledger at any point in time after node provisioning, with cryptographic assurance that no entry has been modified, deleted, or inserted after the fact.}

\subparagraph{Formal statement.}  Let $\mathcal{F}(n)$ denote the Forensic Ledger
entries for node $n$.  Structural Auditability requires that $\mathcal{F}(n)$ forms
a cryptographically chained sequence $f_1, f_2, \ldots, f_k$ in which each
$f_i$ contains $\text{hash}(f_{i-1})$, and that this chain property is
verifiable by any auditor with read access to $\mathcal{F}(n)$ without requiring
access to any other component of the framework.

\subparagraph{Enforcing mechanism.}  The \fact{} layer's Forensic Ledger Writer
assigns each entry a logical timestamp and computes the SHA-256 hash of the
preceding entry before committing.  The chain is initialized with a genesis hash
established at node provisioning.  The \fact{} layer's pre-commit hooks for all
five pillars commit their state transitions and events to the Forensic Ledger
atomically with the state change -- no state transition is committed unless the
corresponding Forensic Ledger entry is also committed.  The cryptographic chaining
is enforced by the \fact{} layer and cannot be bypassed by any machine actor.

\subparagraph{Dependence.}  This property depends on Invariant~\ref{inv:forensic-completeness}
(the completeness of Forensic Ledger recording) and on Invariant~\ref{inv:ledger-immutability}
(the immutability of both ledgers).  Completeness ensures that no relevant event
is omitted; immutability ensures that the recorded events are authentic.

% ---------------------------------------------------------------
\subsection{Property 3: Compulsory Human Escalation}\label{sec:compulsory-escalation}

\theorem[Compulsory Human Escalation]{For any node $n$, if the Bailout Protocol
enters the ESCALATING state at $n$ and all resolution paths within current
authority are exhausted, then the protocol must eventually contact $h(n)$ and
must not enter a terminal resolving state without a directive issued by $h(n)$.}

\subparagraph{Formal statement.}  Let $t_0$ be the time at which the protocol
enters ESCALATING at node $n$.  Let $P(t)$ denote the set of functional
pathways to $h(n)$ available at time $t$.  Let $T_{\max}$ be the deployment-specific
maximum propagation time.  Compulsory Human Escalation requires:
\[
\exists t_1 \geq t_0\ :\ t_1 - t_0 \leq T_{\max}
\ \implies\ h(n) \text{ is contacted}
\]
and, additionally:
\[
\text{state}(n) = \text{RESOLVED}
\ \implies\ \text{directive}(n) \in \{\texttt{RESOLVE},\ \texttt{REJECT}\}
\]
where the directive is verified by the \fact{} layer as issued by $h(n)$.

\subparagraph{Enforcing mechanism.}  The Bailout Protocol's state machine
(Section~\ref{sec:bailout-full}) enforces the compulsory escalation property
through its deterministic transition relation: the only transitions entering
RESOLVED require either a RESOLVE or REJECT directive from $h(n)$.  The
\fact{} layer's pre-commit hook for the Bailout Protocol rejects any directive
whose issuer is not verified as $h(n)$.  The escalation algorithm propagates
via the immutable parent pointer chain, bypassing intermediate machine actors
via the BYPASS property, which is itself enforced by the \fact{} layer's
pre-commit hooks.

\subparagraph{Dependence.}  This property depends on Invariant~\ref{inv:human-anchor}
(the exclusivity of $h(n)$ as terminus), Invariant~\ref{inv:parent-pointer} (the
immutability of the parent pointer chain), and the correct operation of the
Pathway Health Registry ensuring that at least one functional pathway to $h(n)$
is available at all times during normal operation.

% ---------------------------------------------------------------
\subsection{Property 4: Spawn Tree Integrity}\label{sec:spawn-tree-integrity}

\theorem[Spawn Tree Integrity]{The spawn tree of any \bxthree{} system is an
acyclic graph with a designated human root at each node, established at node
provisioning, and the parent pointer of each node is immutable throughout its
operational lifetime.  No node may transfer state to any node outside its
spawn subtree except through the authorized cross-domain channel defined by
the Spatial Firewall pillar.}

\subparagraph{Formal statement.}  Let $\mathcal{T}(n)$ denote the spawn tree
rooted at node $n$.  Spawn Tree Integrity requires: (i) $\mathcal{T}(n)$ is a
directed acyclic graph; (ii) for every node $m \in \mathcal{T}(n)$, the parent
pointer $\alpha(m)$ is established at provisioning and is never modified; (iii)
the human anchor $h(r)$ of the root node $r \in \mathcal{T}(n)$ is a named
human principal; (iv) no state transfer between nodes $p, q \in \mathcal{T}(n)$
occurs except through the authorized channel where $\alpha(p) = q$ or
$\alpha(q) = p$ (parent-child) or where both $p$ and $q$ share a common
ancestor $a$ and the transfer occurs through the Spatial Firewall's authorized
cross-domain channel.

\subparagraph{Enforcing mechanism.}  The \fact{} layer's pre-commit hooks
enforce all four conditions at the pre-commit stage.  The Recursive Spawning
state machine (Section~\ref{sec:pillar-recursive-spawning}) rejects any spawn operation that does
not establish $\alpha(c)$ as the spawning node at provisioning.  The Spatial
Firewall (Section~\ref{sec:pillar-spatial-firewall}) intercepts all cross-domain communication
and permits only the authorized cross-domain channel under authorization ledger
verification.  The Forensic Ledger records all spawn and state-transfer events,
enabling retrospective verification of the acyclicity and isolation properties.

\subparagraph{Dependence.}  This property depends on Invariant~\ref{inv:parent-pointer}
(the immutability of the parent pointer) and on the correct operation of the
Spatial Firewall's pre-commit hook for cross-domain communication.  If either
mechanism is compromised, the spawn tree can no longer be guaranteed acyclic
or isolated.

% ---------------------------------------------------------------
\subsection{Property 5: Sandbox Containment}\label{sec:sandbox-containment}

\theorem[Sandbox Containment]{Every execution event within a sandboxed computation
is contained within the sandbox's provisioned execution budget and isolation
boundary.  No side effect of a sandboxed computation may propagate to the
host system or to any other sandbox without passing through the \fact{} layer's
authorization verification.  If the sandboxed computation attempts to exceed
its execution budget, the abort sequence is triggered and the deviation is
recorded in the Forensic Ledger before any subsequent action is taken.}

\subparagraph{Formal statement.}  Let $s$ be a sandbox instance with execution
budget $B(s)$ and isolation boundary $I(s)$.  Let $\mathcal{E}(s)$ be the set of
execution events generated by the sandboxed computation.  Sandbox Containment
requires: (i) $\forall e \in \mathcal{E}(s)$, $e$ is observable only within $s$
or within the \fact{} layer's monitoring interface; (ii) if the provisioned
execution time $T(s)$ expires before the sandboxed computation completes, the
abort sequence is triggered; (iii) the abort sequence records the termination
event in the Forensic Ledger before the sandbox is destroyed.

\subparagraph{Enforcing mechanism.}  The Sandbox Gate (Section~\ref{sec:pillar-sandbox-gate})
enforces the execution budget and isolation boundary through the \fact{} layer's
pre-commit hook.  Every I/O operation attempted by the sandboxed computation is
intercepted by the gate and submitted to the \fact{} layer for authorization
verification before execution.  The abort sequence is triggered by the watchdog
timer embedded in the Sandbox Gate, which operates with the same scheduling
priority as the Bailout Monitor and is not preemptible by sandboxed code.
The Forensic Ledger records the abort event with the full execution trace
generated during the sandbox session.

\subparagraph{Dependence.}  This property depends on the isolation guarantees
provided by the underlying sandboxing substrate (a deployment assumption) and
on the correct operation of the Sandbox Gate's watchdog timer.  If the watchdog
timer fails to fire or the pre-commit hook is bypassed, the containment
property is violated.
